\section{Qualifications for Initiation}

Note: From time to time, I have been asked about how to go about forming an order. But the desired answer is always about what books to read, or what principles to adhere to. Rather, the purpose of an order is to gain dominion over oneself, and thence over the world. The subtle rules the dense.

\emph{Fear is the only darkness.}

The \textbf{Candidate for Initiation} must demonstrate fearlessness, freedom from lust and finally loyalty and the ability to be silent. The details of these tests were always kept secret. But when outsiders learned bits and pieces of them, they would come to outlandish conclusions. For example, there are many misconceptions about the initiation rites of the \textbf{Templars} based on rumours and partial knowledge.

\paragraph{Facing Danger}


The Candidate must be fearless in the face of unexpected or dangerous things and events from the four elements.

\begin{enumerate}

\item \textbf{Fire:}

The Candidate must pass through fire without fear of being burned. We see presentiments of this in firewalks, for example.

\item \textbf{Water:}

This may involve swimming through rough seas or rivers.

\item \textbf{Air:}

The Candidate stands unsupported on perhaps a square meter platform, while looking down from a great height. Precipices or ship's masts were used. Perhaps nowadays, there are suitable buildings or other man-made structures.

\item \textbf{Earth:}

The Candidate enters a dark and narrow subterranean tunnel or an unknown cave.
\end{enumerate}


\paragraph{Mystical Fear}

This test would include solitary nocturnal visits to cemeteries or haunted houses. We still see this in the practices of the Chodpa sect of Tibetan Buddhism.

\paragraph{Astral Visions}

The Candidate is made to confront terrifying and hideous astral visions, inducing fear and temptations. We see something similar in the story of \textbf{St Anthony of the Desert}, whose visions would frighten neighboring monks away from his cave.

\paragraph{Lust and Passion}

The Candidate must be able to suppress sexual desire, even when its fulfillment is made possible. We see an example of such a test in Miguel Serrano's \textbf{Nos\footnote{\url{https://www.gornahoor.net/?p=70}}}. Even in Sex Magick and Tibetan Tantra\footnote{\url{https://www.gornahoor.net/?p=51}}, the element of lust and pleasure must be separated from any sex act.

\paragraph{Secrecy}

Since the cohesion of any initiatory group requires mutual trust, the Candidate must demonstrate his ability to fulfill a confidential order or instruction, to maintain silence about a secret, and to not betray someone else's personal intention. He must do this of his own free will, as his failure would incur no punishment (other than refusal of initiation). He would be subject to temptations to fail this test.

Further Reading

\textit{Opening the Dragon Gate}, by Chen Kaiguo and Zheng Shunchao

\textit{Shaman of Tibet}, by Heather Hughes-Calero

\textit{Transcendental Magic}, by Eliphas Levi

\textit{The Tarot}, by Mouni Sadhu


\flrightit{Posted on 2010-11-06 by Cologero}

\begin{center}* * *\end{center}

\begin{footnotesize}
\begin{sffamily}
\texttt{Lemonheadz on 2013-01-23 at 12:20 said:}

``The Candidate must be able to suppress sexual desire, even when its fulfillment is made possible.''

Are there techniques to help with this or must a beginner simply not act on his desires?

\hfill

\texttt{Cologero on 2013-01-23 at 23:38 said:}

This is not a topic for a comments section, Lemonheadz. By ordering one's life first, perhaps through prayer, concentration and meditation, the rest will follow. That particular test is not for a ``beginner''; that is why preparation is necessary.

\hfill

\texttt{Lemonheadz on 2013-01-24 at 11:36 said:}

Beautiful.

\hfill

\texttt{Pickman on 2013-08-30 at 23:38 said:}

I get the feeling again with this category of rhetoric that all of the above could be achieved by a cunning business like competitive person if they knew in advance of what the goal to be achieved was. This is why organizations with solid systems get infiltrated and destroyed time and again. Nothing has lasted time, absolutely nothing. So even Gornahoor has not figured that one out. In essence we should be asking ourselves what spiritual tests of men are beyond mere corruption?

\texttt{francismercuri on 2013-08-31 at 17:43 said:}

Pickman: My tendency is to view these test as relatively incorruptible: no amount of ``infiltration'' can alter them, as the ``Trials of the Sphynx'' (and all trials above the Sphynx) are not ``ceremonies''--but, ``real'' internal events, corresponding with the elements, occurring in the Initiate's life, post ``Rite''. Initiation either does establish an actual spiritual link with a super-human ``center'' and hierarchy which ``synergistically'', shall we say, generates living trials that induce the desired changes in an individual seeking Knowledge, preparing him for it--or, Initiation (virtual, actual, or any type we might imagine) is nothing but profane philosophy creatively depicted--and thus, of no real utility.

\hfill

\texttt{Olavsson on 2015-01-27 at 10:28 said:}

Assuming that these tests even correspond to external, physical circumstances, and do not exclusively describe internal events in metaphorical or symbolic language (although I am aware that the tests are meant to inspire inner processes in the candidate for initiation), most of them seem quite straight forward, insofar as the arrangements themselves wouldn't be especially complicated to order.

The trial of Air made me think of Julius Evola's essays about mountaineering, some of which are collected in the book \textit{Meditations on the Peaks: Mountain Climbing as a Metaphor for the Spiritual Quest} (Eng. edition Inner Traditions 1998).

To overcome one's fear of great heights and the risk of injury in the trial of Air, mountain climbing would be a more powerful and effective experience than walking to the top of tall buildings and the like, I am quite certain. But this takes some preparation, as you can't go straight to ascending the Matterhorn without some physical training and experience. This is actually one of my goals, and it would be excellent to be in touch with other traditionalists and spiritual seekers who are also interested in such adventures as a means to test oneself.

When it comes to the Astral Visions test, it gets at once more complicated. How would such experiences be brought about? Hallucinogenics/entheogenics such as DMT? There might be some objections to be raised against that, depending on one's approach to the matter, and I can see a few problems myself, but I don't have any other idea of how vivid ``astral-like'' experiences could be brought about at command for an initiatory test.

I'm even uncertain of how one would go on about arranging the test for mastery of lust and passion. I can easily imagine how one would test oneself in daily, private life to resist problematic temptations of that nature, but how would one arrange such a test for somebody else in the context of a group? Get a woman to volunteer and try to seduce someone?

\hfill

\texttt{Cologero on 2015-01-27 at 10:39 said:}

Interesting thoughts, Olavsson. Re the last test, how about something along these lines: \textit{The Test}\footnote{\url{http://www.gornahoor.net/?p=70}}.

Or, as you point out, it could be understood on a higher level. E.g., this from \textit{Gnosis}, Vol 1, Chapter IX, Section (2):

\begin{quote}\footnotesize
In the same way that on the physical plane the woman bears the fruit of love through her pregnancy, brings it into the world, nourishes it, and educates it, on the

moral plane it is the man who conceives the inspired idea or, fertilized by the woman, develops it inside him, and lastly brings it into the world in the form of some work, or more generally in the form of a creation.
\end{quote}


\hfill

\texttt{Max on 2015-01-27 at 13:18 said:}

I do not see why an initiatory test would have to be separated from the life that we live otherwise. It is not like two separate lives where one is an initiate in one but not in the other, so trials and temptations could naturally take place anywhere and whenever the least expected. Regarding astral visions i think those will also happen quite naturally after a certain awakening to subtler environmental influences have taken place.

Often we hear about initiations in the past where the initiate had first to descend into darkness in order to ascend again, however I do think that there is quite enough darkness around us as it is in this moment of time.

\hfill

\texttt{Olavsson on 2015-01-27 at 14:06 said:}

Max:

I largely agree with what you have to say. Life itself is the grand trial which can lead to `natural' initiations if you only have the eyes to see, the proper understanding and will. The problem is that most of our contemporaries are entirely blind to this dimension of life, unlike the Man of Tradition, whom, as Evola would have said it, saw everything in this world as a symbol reflecting something profound -- or, if you will, as shadows or reflections of higher eternal principles. One internal change that must accompany the person's liberation from the prison cell that is the modern mind, is to re-establish that higher awareness of what life truly is in a higher sense -- to see the subtle forces and principles that permeate the world and which are at work within one's own being. You open your eyes to new, subtler dimensions of life that are ignored, suppressed by the gross material and mechanical outlook typical of the modern mind. If one, as a modern man, does not actively seek to tune onto such currents and attain experience of subtler states of being -- a sort of awareness that is necessary for initiations to occur -- then one will in all likelihood go through all of life unaware of these things. Hence most people will not naturally experience initiations if they do not make a decision to actively aim for it, like an arrow reaching its target. While the kind of tests or trials that Cologero describes in his post do not in themselves constitute THE thing, they represent a conscious and deliberate attempt at dramatization, which, if experienced with sincerity and right intention can help bring about at least a small shift in consciousness and awareness, making it easier for the aspiring modern candidate for initiation to enter a path and look at the world with new eyes.

\hfill

\texttt{Max on 2015-01-28 at 13:51 said:}

In Hrolf Kraki's Saga, the god is disguising himself as a farmer and offers the traveling king a sword, a shield and an armour that the king turns down.

``But they had gone barely a mile, enough for the mist to lift in their heads, when Bjarki stopped. The rest did likewise. Dim in twilight, he told them: `Too late do the unwise come to understanding. So it is with me. I have a feeling we did not behave very sagely when we said no to that we ought to have said yes to. We may have bidden victory go from us.' ''

Saying no is at least better than the apathy of not using the will at all, and the good news is that a recognized mistake is redeemable as well as a learning experience.

\hfill

\texttt{Max on 2015-01-28 at 14:49 said:}

Actually that was a badly chosen translation in my comment as it is from a retelling of the story. This one should be more true to the source.

``When they had not gone very far, Bodvar halted and said, `Good sense comes late to fools, and so it comes to me now. I suspect that we have not behaved very wisely in rejecting what we should have accepted. We may have denied ourselves victory.'\,''

\hfill

\end{sffamily}
\end{footnotesize}

